
Reliable grounded VideoQA depends not only on predicting the correct answer, but also on identifying the temporal evidence that truly supports it. This requirement is central to grounded video understanding because answer correctness alone does not guarantee that a prediction is visually substantiated. \emph{NExT-GQA} shows that strong answer accuracy can still arise from weak evidence alignment \cite{xiao2024can}, while \emph{ReXTime} further demonstrates that the required evidence may be temporally displaced from the question context, making grounded reasoning substantially harder \cite{chen2024rextime}. These observations suggest that reliable grounded VideoQA depends not only on reasoning ability, but also on how effectively answer-relevant evidence is acquired and verified.

Recent Multimodal Large Language Models (MLLMs) have substantially improved video understanding \cite{tang2025video, wu2025survey_vtg_mllm}. In temporal grounding, \emph{VTimeLLM} enhances timestamp-sensitive modeling \cite{huang2024vtimellm}, \emph{TimeChat} strengthens long-video temporal perception \cite{ren2024timechat}, and \emph{Grounded-VideoLLM} improves fine-grained grounding through grounding-oriented training and architectural design \cite{wang2025grounded_videollm}. In video reasoning, \emph{SeViLA} decomposes answering into localization and language reasoning \cite{yu2023sevila}, \emph{LLoVi} studies efficient long-range VideoQA with modular processing \cite{zhang2024simple}, and \emph{VideoMind} introduces a role-based workflow with specialized planner, grounder, verifier, and answerer modules \cite{liu2025videomind}. These methods show that temporal modeling and role specialization are effective design principles. However, they still leave a key challenge underexplored in grounded multiple-choice VideoQA. The language available for answer prediction is often not the language needed for temporal evidence retrieval.

The central hypothesis of this work is that grounded VideoQA benefits from explicitly separating \emph{answer discrimination} from \emph{evidence retrieval}. In multiple-choice VideoQA, answer candidates are written to distinguish options, not to describe visually localizable events. They are therefore often abstract, incomplete, or weakly visual. By contrast, temporal grounding requires retrieval queries that expose observable evidence. Directly grounding raw answer candidates can therefore produce segments that are temporally nearby but semantically weak, noisy, or only partially aligned with the evidence needed for answering. This mismatch is particularly harmful on benchmarks such as \emph{NExT-GQA} and \emph{ReXTime}, where prediction quality is evaluated together with evidence quality \cite{xiao2024can, chen2024rextime}.

This paper presents \emph{EviR}, a modular framework for temporally grounded VideoQA that separates \emph{evidence reformulation}, \emph{evidence acquisition}, \emph{evidence validation}, and \emph{answer generation}. In the current multiple-choice instantiation, a Refiner rewrites each answer candidate into a grounding-friendly query, a Retriever localizes candidate temporal spans, a Validator scores and ranks the retrieved candidates by evidence reliability, and a Generator predicts the final answer from the highest-ranked evidence clip. This design follows the view that grounded answer prediction should be built on explicit evidence quality control rather than on direct end-to-end generation from the full video.

Evaluation on \emph{NExT-GQA} and \emph{ReXTime} shows that this design consistently improves strict grounding-aware metrics. EviR improves IoU@0.5 by 5.77 over the strongest prior baseline on \emph{NExT-GQA} and improves Acc@IoU=0.5 by 7.05 on \emph{ReXTime}, while also achieving the best or comparable answer accuracy. These results support the hypothesis that better alignment between answer formulation and temporal retrieval leads to more reliable grounded reasoning.

The main contributions are summarized below.
\begin{itemize}
    \item A key mismatch in grounded VideoQA is identified between language formulation and temporal evidence retrieval. On current multiple-choice benchmarks, answer candidates are useful for discrimination but often unsuitable for evidence localization.
    \item \emph{EviR} is presented as a modular framework that resolves this mismatch through evidence-guided query reformulation, temporal retrieval, explicit evidence validation, and grounded answer generation.
    \item Strong results are achieved on \emph{NExT-GQA} and \emph{ReXTime}, showing that improving evidence quality leads to more accurate and more reliable grounded VideoQA under strict grounding-aware evaluation.
\end{itemize}

The remainder of this paper is organized as follows. Section~2 reviews related work on temporal grounding, grounded VideoQA, and modular video reasoning. Section~3 presents the EviR framework and its components. Section~4 describes the experimental setup and reports the main results and ablations. Section~5 concludes the paper and discusses future directions.
